% !PRINTED 39
\weberpage{39}
But if $f(x_0)$ vanishes, the development of $R(x)$ takes the form:
\weq{6}{R(x) = A_m(x-x_0)^{m-1/2} + A_{m+1}(x-x_0)^{m-1/2+1} + A_{m+2}(x-x_0)^{m-1/2+2} + \cdots,}
and if $f(x)$ is of the third degree, then in the neighbourhood of the infinitely distant point the development holds
\weq{7}{R(x) = C_m x^{-m-3/2} + C_{m+1} x^{-m-5/2} + C_{m+2} x^{-m-7/2} + \cdots}
From the developments (4) to (7) we can derive the developments of $\Omega$, an additive constant remaining undetermined.

The development of $\Omega$ is then likewise a power series, to which, if $m$ is negative, there is added in the cases (4) or (5) a further term $A_{-1}\log(x-x_0)$, $C_{-1}\log x$.

\gesperrt{We now call $x_0$ a point of the first kind of} $\Omega$ if $m$ is not negative in these developments.

\gesperrt{Then $\Omega$ is finite at the point $x_0$ (also in the case $x_0 = \infty$).}

\gesperrt{The point $x_0$ is called a point of the second kind of $\Omega$,} if in the developments (4), (5) $m$ is negative and $A_{-1} = 0$, $C_{-1} = 0$, and in the case of the development (6), (7) with negative $m$.

\gesperrt{Finally, $x_0$ is called a point of the third kind,} if in the developments (4) or (5) $A_{-1}$ or $C_{-1}$ is different from zero, that is, if a logarithmic term occurs in the development of $\Omega$.

In the integral (2) the zero points of $f(x)$ cannot be of the third kind; they can, however, become of the third kind if we consider the elliptic integral in the general form $\int \Phi(x, \sqrt{f(x)})\,dx$.

The integral $\Omega$ is called an \gesperrt{integral of the first kind} if it has only points of the first kind. Such an integral then has a finite value for all finite and infinite values of $x$.

$\Omega$ is called an \gesperrt{integral of the second kind} if it has only points of the first and second kind, and

$\Omega$ is called an \gesperrt{integral of the third kind} if, besides points of the first and second kind, it also has points of the third kind.
